\section{Introduction}\label{sec:introduction}

For a semilinear initial-value problem
\begin{equation}
  y'(t)=Ly(t)+N(t,y(t)),
  \label{eq:semilinear-intro}
\end{equation}
exponential integrators incorporate the linear flow generated by $L$ into the
time-discretization coefficients.  This idea supports high-order methods for
stiff and highly oscillatory problems while replacing repeated linear solves by
actions of the exponential and the related $\varphi$
functions~\cite{HochbruckOstermann2010}.  The numerical method and the
algorithm used to evaluate those actions are conceptually distinct, but
software often entangles them.  That coupling becomes increasingly restrictive
when the same time integrator should be used with a diagonal spectrum, a dense
reference matrix, a large sparse operator or a matrix-free linear map.

EXPINT was an early software response to this problem.  Its MATLAB catalogue
made exponential Runge--Kutta, multistep and general linear integrators
available in a common package and emphasized rapid deployment and
testing~\cite{BerlandSkaflestadWright2007}.  Since then, matrix-function-vector
algorithms have evolved substantially.  Krylov algorithms for combinations of
$\varphi$ functions, from Expokit to Algorithm~919 and
KIOPS~\cite{Sidje1998,NiesenWright2012,GaudreaultRainwaterTokman2018}, avoid
materializing matrix functions.  Polynomial methods based on Leja
interpolation provide a different tradeoff when usable spectral information is
available~\cite{CaliariVianelloBergamaschi2004,CaliariKandolfOstermannRainer2016,DekaEinkemmerTokman2023}.
General-purpose scientific libraries also provide scalable exponential actions
such as \texttt{expm\_multiply}~\cite{AlMohyHigham2011,VirtanenEtAl2020}.  These
developments motivate a software boundary in which a time integrator requests
matrix-function actions without being rewritten for each numerical
realization.  The boundary itself is not new: Expokit and Algorithm~919 are
reusable action routines, and OrdinaryDiffEq.jl switches between cached matrix
functions and Krylov actions with one option.  What is new here is what is
built on it.

\pkg{} is an independent Python framework built around that boundary.  Its
design is not a literal translation of the historical EXPINT source.  The
historical catalogue is reconstructed as mathematical data in a single
\texttt{MethodSpec} representation; the execution engine interprets those data;
and action backends implement the numerical realization.  This separation
allows the same method specification to operate on dense, diagonal, sparse and
matrix-free problems and makes it possible to instrument work at the
matrix-function layer.  It also exposes a useful verification surface: method
metadata, coefficient identities, history layout, action algebra and backend
statistics can be tested separately rather than being hidden inside a
method-specific solver.

The software contribution is therefore broader than ``EXPINT in Python'' but
narrower than a claim to a new universally superior exponential solver.  The
article makes five main contributions.
\begin{enumerate}[leftmargin=*,itemsep=2pt]
  \item We encode a 47-entry historical catalogue in a common explicit exponential
  general linear method (EGLM) representation, with separate
  metadata for classical, strong stiff and weak or conditional order.
  \item We introduce an analytic operator-expression layer and a backend
  protocol that decouple method coefficients from matrix-function-action
  algorithms.  Compatible $\varphi$ actions can be fused automatically,
  including sums with distinct right-hand sides when the backend supports
  them, and external backends enter through one hook.
  \item We provide fixed-step execution with explicit startup policies for
  history methods, generic step-doubling adaptivity and the embedded ETD34
  pair with action-tolerance coupling, failure rejection and explicit
  termination rules, and we document a limitation of the published ETD34
  estimator, which vanishes for state-independent forcing.
  \item We combine structural and coefficient checks with a numerical
  cross-check against the original EXPINT implementation.  The supplied
  summary reports near-roundoff agreement for 42 of the 47 entries,
  agreement with a one-line correction for four entries, and regenerated
  output coefficients for one entry.  The summary and its limitations are
  made explicit in Section~\ref{sec:crosscheck}.
  \item We supply work instrumentation and reproducibility material that
  separates portable quantities, namely error, matrix-vector products,
  nonlinear evaluations, accepted and rejected steps and Krylov or polynomial
  diagnostics, from machine-dependent elapsed time.  Historical runs, author-reported candidate runs and manuscript
  postprocessing are distinguished in the experimental record.  Timing
  results are retained as paired observations.  A broadband problem without
  manufactured forcing is also included, with its original numerical record, reference state and individual timing
  repetitions retained in the private companion.
\end{enumerate}

This scope determines what the article does not claim.  No runtime
superiority over external packages such as EXPODE, LeXInt, \texttt{rkstiff} or
the reference KIOPS implementation is asserted; the relationship to them is
functional (Section~\ref{sec:related}).  No distributed-memory scalability
claim is made because the release has no distributed backend interface.  The
rule-based automatic backend is presented as configurable and
problem-dependent, not as a globally optimal selector; on the canonical
timing problem it chooses the slower backend.  And the default numerical
startup does not in general preserve the nominal orders six and seven of the
high-order history methods.

The remainder of the paper follows the software boundary itself.
Section~\ref{sec:math} states the common mathematical representation and
action algebra.  Section~\ref{sec:architecture} describes external state,
startup, execution, work accounting and the release-facing API.
Section~\ref{sec:catalogue} presents the validated historical catalogue, its
reconciliations and the cross-check against the original sources.
Section~\ref{sec:backends} describes the action backends and fusion contract,
and Section~\ref{sec:adaptivity} the adaptive interfaces.  The numerical
program is reported in Section~\ref{sec:experiments}.  Related software is
positioned in Section~\ref{sec:related}; reproducibility, AI-assisted research
use and limitations are discussed in Section~\ref{sec:repro}; and
Section~\ref{sec:conclusions} concludes.
